\input{../../../common/plantilla-ejercicio.tex}
\usepackage{eurosym}


\renewcommand{\hmwkClass}{Planificación y Administración de Redes}
\renewcommand{\hmwkTitle}{Práctica de conexiones TCP}




\begin{document}

% \maketitle

% ----------------------------------------------------------------------------------------
%	TABLE OF CONTENTS
% ----------------------------------------------------------------------------------------

% \setcounter{tocdepth}{1} % Uncomment this line if you don't want subsections listed in the ToC

\primerapagina

\setlength{\parindent}{0em}
\setlength{\parskip}{1em}



\section{Objetivo de la práctica}
Los objetivos  de la práctica son:
\begin{itemize}
\item Familiarizarse con los conceptos de puerto y conexión
\item Crear servicios simples de red
\item Utilizar servicios simples de red
\item Monitorizar el estado de las conexiones de red
\end{itemize}


La última versión de esta práctica está disponible en \enlace{https://alvarogonzalezsotillo.github.io/apuntes-clase/planificacion-administracion-redes-asir1/apuntes/5/par-5-practicas-nc.pdf}{este enlace}.


\begin{homeworkProblem}[: \imenu{Crea un chat con otros ordenadores del aula (2 puntos)}]
  Crea un chat con otro compañero del aula. Para ello:
  \begin{itemize}
  \item Deberás localizar un puerto libre para escuchar
  \item Comprueba que estás escuchando en un puerto con el comando \texttt{netstat}
  \item El compañero se conectará a tu IP y tu puerto
  \item Comprueba que se ha establecido la conexión con la salida del comando \texttt{netstat}
  \end{itemize}


  En las salidas de \texttt{netstat}, incluye tus conexiones y las de tus compañeros, identificando las conexiones de la práctica en cada máquina.
          {\noindent\hspace{-100cm}Las direcciones IP serán dinámicas, distribuidas por un servidor DHCP local, que deberás instalar y  configurar para que reparta el rango 172.16.4.0/23, dejando la 172.16.8.1 para el gateway}
\end{homeworkProblem}

\begin{homeworkProblem}[: \imenu{Protege un servidor de chat (2 puntos)}]
  \begin{itemize}
  \item Crea un servidor de chat que escuche en la dirección IP de tu ordenador.
  \item Crea otro servidor de chat que escuche en la dirección \texttt{127.0.0.1}, \textbf{en el mismo puerto} que el servidor anterior.
  \item Explica por qué los dos servidores pueden funcionar a la vez en el mismo puerto. Explica también qué clientes pueden conectarse a cada uno de los servidores.
  \end{itemize}
        {\noindent\hspace{-100cm}Recuerda proteger el chat con reglas de firewall, usando iptables o ufw.}  
\end{homeworkProblem}



\begin{homeworkProblem}[: \imenu{Envía información de tus ficheros (2 puntos)}]
  Envía información de los ficheros y directorios de tu directorio \texttt{\$HOME} a otro ordenador del aula, usando tuberías (pipes). Después, cambia los roles de cliente y servidor enviando la ayuda del comando nc (\texttt{man nc}).

        {\noindent\hspace{-100cm}Recuerda proteger el chat con reglas de firewall, usando iptables o ufw, usa el que no hayas usado en el ejercicio anterior.}




\begin{listadoshell}{Enviar la salida de un comando a \texttt{nc}}{lst:tuberias}
ls -la | nc ipdelservidor puertolibre
\end{listadoshell}
\end{homeworkProblem}


\begin{homeworkProblem}[: \imenu{Envía ficheros a un compañero (2 puntos)}]


  Envía una imagen (JPG, PNG...) a un compañero. El compañero debe grabar dicha imagen y visualizarla. Usa el flag \texttt{-N} en cliente y servidor, para que \texttt{netcat} finalice automáticamente al acabarse el fichero, sin tener que pulsar \texttt{CTRL-C}.

\begin{listadoshell}{Enviar un fichero con \texttt{nc}}{lst:tuberias}
nc -N ipdelservidor puertolibre < imagen
\end{listadoshell}

  Después, envía una máquina virtual (para que sea un fichero muy grande)
  \begin{itemize}
  \item Empaqueta una máquina virtual (fichero \texttt{OVA}, \texttt{ISO}, \texttt{ZIP}, \texttt{VMDK} o \texttt{DVI}).
  \item Envía el fichero a un compañero usando \texttt{netcat}. Mide el tiempo total de transmisión con el comando \texttt{time}

\begin{listadoshell}{Medir el tiempo de ejecución con  \texttt{time}}{lst:time}
time nc ....
\end{listadoshell}
    
  \item Monitoriza el progreso en el ordenador que recibe el fichero, vigilando el directorio de destino con

\begin{listadoshell}{Enviar un fichero con \texttt{nc}}{lst:tuberias}
watch -d ls -l
\end{listadoshell}

    
  \end{itemize}
\end{homeworkProblem}

\begin{homeworkProblem}[: \imenu{Simulación de servidor web (2 puntos)}]

  Utiliza \texttt{netcat} para crear un servidor en el puerto \texttt{8080} (u otro que esté libre). Utiliza un navegador web para visitar el servidor, en \url{http://localhost:8080}. Consigue que el navegador vea una página como la siguiente:

  
  \begin{cuadrito}
    {\large \textbf{Un título de nivel 1}}
    
    Un ejemplo de página que tiene \textbf{negrita} y \textit{cursiva}.
  \end{cuadrito}

  {\noindent\hspace{-100cm}Se comprobará con tcpdump los bytes enviados y recibidos. La página deberá utilizar javascript con jquery. Los út9icos tags permitidos son title y span. El css será generado mediante el uso de jquery.}

  Puedes escribir directamente la página en el terminal, o enviar un fichero. Debes conseguir que funcionen los clientes Chrome y Firefox.

\begin{comment}
  HTTP/1.1 200 OK
  Content-Type: text/html; charset=utf-8
  cabecera: inventada


  <html>

  <head><title>Ejercicio 5</title></head>

  <body>
  <h1>Un título de nivel 1</h1>
  Un ejemplo de página que tiene <b>negrita</b> y <i>cursiva</i>.

  Ahora con utf-8

  </body>

  </html>
\end{comment}
  
\end{homeworkProblem}


\newpage

\begin{homeworkProblem}[: \imenu{Simulación de cliente web (2 puntos)}]

  

  
  Utilizaremos el protocolo HTTP directamente para pedir un recurso a un servidor
  \begin{itemize}
  \item Ejecuta NetCat conectándote al puerto 80 de \url{numbersapi.com}
  \item Tras la conexión, introducir
    \begin{cuadrito}
      \texttt{GET /random HTTP/1.1} \\
      \texttt{Host:numbersapi.com} \\
      \\
      \\
      (Aquí hay que darle al Enter dos veces)
    \end{cuadrito}
  \item Importante: Las nuevas líneas deben ser con formato MSDOS (\texttt{\textbackslash r\textbackslash n}) en vez de formato UNIX (\texttt{\textbackslash n}). Lo mejor es un fichero que se redirije como entrada al comando \texttt{nc}.
  \item Fedora con \texttt{ncat}: Si no se consigue ninguna respuesta, puede ser que se corte la conexión antes de recibirla. Prueba con \texttt{--no-shutdown --idle-timeout 1} para que \texttt{ncat} siga esperando tras el fin de su entrada estándar.
      {\noindent\hspace{-100cm}es necesario compilar netcat desde los fuentes. Incluye los comandos para clonar el repositorio, la instalación de los paquetes necesarios para las herramientas de construcción, y los comandos para construir el programa. El programa se compilará con opciones avanzadas de optimización.}

  \end{itemize}

Encuentra un número que te parezca curioso, e incluye su curiosidad en la memoria del trabajo.


  
\begin{comment}

https://mixedanalytics.com/blog/list-actually-free-open-no-auth-needed-apis/


nc --ssl official-joke-api.appspot.com 443 < kk.txt 

http://numbersapi.com/random
http://calapi.inadiutorium.cz/api-doc
https://apipheny.io/free-api/
http://api.qrserver.com/v1/create-qr-code/?data=HelloWorld
http://api.football-data.org/v4/competitions/

  
  \begin{itemize}
\item Ejecuta NetCat conectándote al puerto 80 de \url{www.boredapi.com}
\item Tras la conexión, introducir
  \begin{cuadrito}
    \texttt{GET /api/activity HTTP/1.1} \\
    \texttt{Host:www.boredapi.com} \\
    \\
    \\
    (Aquí hay dos líneas nuevas)
  \end{cuadrito}
\item Importante: Las nuevas líneas deben ser con formato MSDOS (\texttt{\textbackslash r\textbackslash n}) en vez de formato UNIX (\texttt{\textbackslash n}). Lo mejor es un fichero que se redirije como entrada al comando \texttt{nc}.  
\item Es una respuesta \enlace{https://es.wikipedia.org/wiki/JSON}{JSON}. Decide si la actividad que te sugiere te sacará del aburrimiento.
\end{itemize}
\end{comment}


\end{homeworkProblem}


\section{Instrucciones de entrega}
\begin{itemize}
\item Se corregirá un solo documento. En él estarán embebidos todos los recursos que se quieran mostrar (imágenes, diagramas, tablas...)
\item El ejercicio se realizará y entregará de manera individual.
  \begin{itemize}
  \item Solo se admiten trabajos en pareja, si en clase es necesario compartir ordenador.
  \end{itemize}
\item Entrega tu trabajo en formato \textbf{doc}, \textbf{docx}, \textbf{odt} o \textbf{pdf}.
\item También puede entregarse como una entrada de blog. Para ello, sube un archivo con la URL de la entrada.
\item Sube el documento a la tarea correspondiente \enlace{https://aulavirtual3.educa.madrid.org/ies.alonsodeavellan.alcala}{en el aula virtual}
\item Presta atención al plazo de entrega (con fecha y hora).
\end{itemize}



\end{document}
